\section{Kết quả, và ranh giới của nó}
\label{sec:ch12-ket-qua-va-ranh-gioi}

\index{chỉ số phản thực!bảy chỉ số}%
Mục trước dựng xong hai vế của phép đối chiếu ở phía người. Vế còn lại là bảy
chỉ số tự động, tính trên đúng 85 lời giải thích ấy. Bài báo lấy chúng từ tài
liệu về lời giải thích phản thực và cho mỗi chỉ số một công
thức~\autocite{p25userperception}. Sách mô tả từng chỉ số bằng lời và không
chép công thức nào, vì cái chương này lập luận là quan hệ giữa bảy con số ấy
với đánh giá của người, và quan hệ đó không đọc được từ công thức. Đây là lựa
chọn của sách, không phải giới hạn của nguồn: khác với bảng tính chất mà
chương~\ref{ch:lime-nao-dang-tin} đọc, vốn không kèm thủ tục tính nào, bài báo
này có đủ công thức cho cả bảy.

Sparsity đếm số đặc trưng khác nhau giữa đầu vào gốc và bản đã sửa, nên càng
thấp thì lời giải thích càng đổi ít chỗ. Proximity đo khoảng cách giữa hai
điểm ấy trong không gian đặc trưng, nên càng thấp thì bản đã sửa càng sát bản
gốc. Closeness đo khoảng cách trung bình từ bản đã sửa tới những điểm huấn
luyện gần nó nhất, nên càng thấp thì bản đã sửa càng nằm gần dữ liệu đã quan
sát được. Diversity đo mức độc lập lẫn nhau giữa các đặc trưng bị đổi: với mỗi
cặp đặc trưng bị đổi, bài báo tính thông tin tương hỗ chuẩn hóa của cặp ấy,
tức mức biết giá trị của đặc trưng này thì đoán được giá trị của đặc trưng kia,
quy về thang từ 0 tới 1, rồi lấy trung bình phần bù trên mọi cặp; nên càng cao
thì các chỗ đổi càng ít trùng lặp thông tin. Oracle Score nhân xác suất mà mô
hình nền gán cho lớp đích với xác suất mà một mô hình thứ hai, huấn luyện độc
lập, gán cho cùng lớp ấy, nên càng cao thì hai mô hình càng đồng ý rằng bản đã
sửa thuộc lớp đích. Trust Score cần một quy ước về khoảng cách từ một điểm tới
một lớp, và quy ước ấy là khoảng cách trung bình tới ba điểm huấn luyện gần
nhất thuộc lớp đó; chỉ số lấy tỉ số giữa khoảng cách tới lớp gần nhất khác và
khoảng cách tới chính lớp được dự đoán, nên càng cao thì bản đã sửa càng nằm
sâu trong lớp của nó. Completeness hỏi lời giải thích có động vào những đặc
trưng mà mô hình coi trọng hay không: lấy năm đặc trưng có \tn{mức quan
trọng}{importance} SHAP lớn nhất tại điểm gốc, cộng các mức quan trọng ấy lại,
rồi xem phần đóng góp của những đặc trưng vừa nằm trong năm đặc trưng đó vừa bị
lời giải thích đổi chiếm bao nhiêu phần của tổng; nên càng cao thì lời giải
thích càng động vào đúng những chỗ mô hình coi
trọng~\autocite{p25userperception}.

Hai chỗ trong bảy chỉ số ấy dùng lại chữ mà sách đã cấp cho thứ khác, và cả
hai phải nêu ra ở đây. Sparsity của chương~\ref{ch:khung-hop-nhat} là đại
lượng đếm cần bao nhiêu phần tử điểm cao mới đạt fidelity, tức một tính chất
của lời giải thích attribution; sparsity ở đây đếm số đặc trưng mà một lời
giải thích phản thực đổi. Còn \tn{tính đầy đủ}{completeness} của
chương~\ref{ch:attribution-theo-gradient} là tiên đề buộc tổng attribution
bằng chênh lệch đầu ra giữa đầu vào và baseline; Completeness ở đây là một tỉ
lệ khối lượng SHAP và không phải tiên đề nào cả. Sách giữ cả bảy cái tên bằng
tiếng Anh, đúng theo lý do mà chương~\ref{ch:do-do-trung-thuc} đã nêu cho các
tên chỉ số của nó, và ghi hai chỗ va chữ ở đây một lần.

\index{chỉ số phản thực!dựng trên SHAP}%
Chỉ số cuối cùng trong bảy đáng một câu riêng vì nó khép một vòng mà cả
Phần~IV đang mở. Completeness được tính bằng SHAP, tức bằng chính phương pháp
mà chương~\ref{ch:shap} đọc và chương~\ref{ch:do-do-trung-thuc} chưa xác nhận
được độ trung thực. Vậy một trong bảy con số dùng để chấm chất lượng lời giải
thích lại là hàm của một lời giải thích khác, và nếu attribution kia sai thì
con số này sai theo mà không có dấu hiệu nào trên mặt số. Quan sát này là của
sách; bài báo chỉ nêu SHAP như cách tính, không bàn tới hệ quả ấy.

\index{chỉ số phản thực!tương quan với đánh giá của người}%
Kết quả thứ nhất là tương quan từng cặp. Các tác giả tính hệ số tương quan
Pearson giữa bảy chỉ số và năm chiều đánh giá cộng với điểm gộp, tính riêng
cho từng bộ dữ liệu. Gộp toàn bộ 85 lời giải thích lại thì chỉ một chỉ số duy
nhất có liên hệ đạt mức ý nghĩa thống kê với điểm gộp: Trust Score, với
$r = 0.307$ và $p = 0.004$. Sáu chỉ số còn lại đều cho tương quan không đáng
kể, trị tuyệt đối dưới 0.1~\autocite{p25userperception}.

Kết quả thứ hai nằm dưới kết quả thứ nhất và nặng hơn nó. Khi tách theo bộ dữ
liệu, các tương quan không chỉ yếu mà còn ngược chiều nhau. Trên Mushroom,
sparsity, diversity, proximity và closeness tương quan âm ở mức vừa tới mạnh
với chiều đủ chi tiết, chiều hài lòng và điểm gộp, hệ số từ $-0.38$ tới
$-0.64$; bài báo đọc điều đó là người dùng ở lĩnh vực này thích những lời giải
thích đổi ít đặc trưng và đổi ở mức nhỏ. Trên Obesity Levels thì phần lớn tương quan mang
dấu dương: diversity, Trust Score và Completeness tương quan dương với chiều
hợp lý, chiều hài lòng và điểm gộp, hệ số từ 0.37 tới 0.52, và bài báo đọc
điều đó là người dùng ở lĩnh vực này thiên về những lời giải thích đầy đủ
thông tin hơn. Trên Heart Disease thì tương quan yếu đều và không đạt mức ý
nghĩa ở mọi cặp, dấu thì lẫn lộn. Độ lệch chuẩn trung bình của các hệ số tương
quan giữa ba bộ dữ liệu là 0.31~\autocite{p25userperception}.

Con số 0.31 ấy là con số quan trọng nhất trong cả mục này, quan trọng hơn
$r = 0.307$ ở trên, và lý do thì nằm ở chỗ một dụng cụ đo được dùng để làm gì.
Một chỉ số tương quan yếu với \tn{cảm nhận của người dùng}{user perception}
vẫn còn dùng được nếu nó yếu
đều: ta biết nó đo lệch bao nhiêu và hiệu chỉnh. Một chỉ số đổi cả dấu khi
chuyển sang bộ dữ liệu khác thì không hiệu chỉnh được, vì muốn hiệu chỉnh phải
biết trước dấu, mà biết trước dấu thì đã phải chạy nghiên cứu người dùng trên
chính bộ dữ liệu ấy rồi, tức đã không cần tới chỉ số nữa. Bài báo tự phát biểu
điều đó ngay sau khi in con số: không chỉ số đơn lẻ nào, và cũng không tổ hợp
chỉ số nào giữ được ổn định, làm đại lượng thay thế đáng tin cho đánh giá của
người xuyên qua các lĩnh vực~\autocite{p25userperception}.

\index{chỉ số phản thực!tổ hợp nhiều chỉ số}%
Vế thứ hai của câu vừa dẫn, về tổ hợp, là kết quả thứ ba và nó được kiểm bằng
một phép thử vét cạn. Trong thực hành, người ta hiếm khi báo cáo một chỉ số
đơn lẻ; người ta in cả một hàng chỉ số cạnh nhau và ngầm cho rằng cả hàng ấy
gộp lại thì mô tả được chất lượng. Các tác giả đem giả định ấy ra thử: họ chạy
mọi tổ hợp con khác rỗng của bảy chỉ số, tức 127 tổ hợp, huấn luyện năm lớp mô
hình dự báo trên mỗi tổ hợp, cho ba bộ dữ liệu và sáu biến mục tiêu, tức năm
chiều đánh giá cộng với điểm gộp, rồi đo bằng $R^2$ tính theo lối đánh giá
chéo năm phần: chia dữ liệu làm năm phần, lần lượt lấy bốn phần huấn luyện và
phần còn lại để chấm, rồi lấy trung bình năm lần chấm ấy, nên con số báo cáo
là mức khái quát ra dữ liệu chưa thấy chứ không phải mức khớp trên dữ liệu
huấn luyện~\autocite{p25userperception}.

Cách đo ấy cần một câu giải nghĩa vì dấu của nó mang thông tin. Hệ số xác định
$R^2$ so sai số của mô hình dự báo với sai số của một mốc tầm thường là luôn
đoán giá trị trung bình. Bằng 0 nghĩa là mô hình không hơn gì cái mốc ấy; âm
nghĩa là nó tệ hơn cả việc đoán trung bình.

Với cách đọc đó, kết quả gọn lại thành một dòng: trung bình trên mọi thiết
lập, cả năm lớp mô hình đều cho $R^2$ âm. Hồi quy tuyến tính trung bình
$-1.253$, XGBoost còn tệ hơn ở $-1.874$ do quá khớp trên cỡ mẫu nhỏ, k láng
giềng gần nhất $-0.887$, và rừng ngẫu nhiên tốt nhất trong năm nhưng vẫn
$-0.474$; mô hình cộng tính tổng quát, lớp mô hình cho mỗi biến dự báo một
đường cong riêng rồi cộng các đường cong lại, thì thường không hội tụ trên cỡ
mẫu này nên không có con số~\autocite{p25userperception}. Ở trường hợp thuận lợi nhất mà bài báo đem ra
minh họa, dự báo chiều hài lòng trên Heart Disease, rừng ngẫu nhiên đạt $R^2$
dương ở 95 trong 127 tổ hợp, nhưng trung bình chỉ 0.067 và cao nhất 0.331.

Còn một điều mà phép vét cạn này thấy được và một phép thử trên vài tổ hợp thì
không: thêm chỉ số vào không những không giúp mà còn làm hỏng. Xét theo số chỉ
số có mặt trong tổ hợp, hồi quy tuyến tính khá lên chút ít tới mức bốn chỉ số
rồi tụt mạnh; rừng ngẫu nhiên đạt đỉnh ở ba và bốn chỉ số, cao nhất $R^2$ bằng
0.33 với ba chỉ số, rồi giảm đơn điệu khi thêm nữa. Bài báo tóm lại rằng ngay
với mô hình phi tuyến, tăng số chỉ số quá khoảng ba tới bốn thì không cải
thiện mà làm giảm hiệu năng~\autocite{p25userperception}.

Đọc theo chiều của một dụng cụ đo, kết quả ấy nói một điều cụ thể. Nếu bảy chỉ
số cùng đo bảy mặt khác nhau của một đại lượng chung, thì đưa thêm mặt vào
phải làm bức tranh đầy hơn. Việc thêm chỉ số làm hỏng dự báo có nghĩa là những
chỉ số thêm vào mang nhiễu chứ không mang thông tin bổ sung về cái mà người
đọc cảm nhận. Bài báo gọi đó là một giới hạn mang tính cấu trúc trong cách các
chỉ số hiện nay nắm bắt những tiêu chí có ý nghĩa với con
người~\autocite{p25userperception}.

\index{độ trung thực!của attribution đặc trưng, phép thử đồng thuận}%
Cái không đi qua được biên giới sang các họ chỉ số của
chương~\ref{ch:do-do-trung-thuc} là mọi con số ở trên. Hệ số 0.307, độ
lệch 0.31, mọi giá trị $R^2$: tất cả đều là số đo trên bảy chỉ số của lời giải
thích phản thực, trên ba bộ dữ liệu bảng, với một phương pháp sinh duy nhất.
Không con số nào trong đó nói gì về deletion, về comprehensiveness, về ROAR
hay về Sensitivity-n. Đây là cùng một ranh giới mà
mục~\ref{sec:ch09-doi-tuong-khac-nhau} dựng cho bài báo neo, và nó được giữ ở
đây vì đúng cùng một lý do.

Cái đi qua được thì là câu hỏi chứ không phải con số: các chỉ số
tự động của một họ lời giải thích đã bị đem đối chiếu với cảm nhận của người
đúng một lần, và chúng không trụ được. Với các họ chỉ số của
chương~\ref{ch:do-do-trung-thuc}, phép đối chiếu ấy chưa được làm lần nào. Cái
đi qua biên giới vì thế là một tiền lệ: phép thử ấy khả thi, và kết quả của nó
không hiển nhiên.

Đi cùng tiền lệ ấy là một mẩu bằng chứng mà bài báo này thuật lại chứ không tự
làm, và nó chạm đúng vào các họ chỉ số của
chương~\ref{ch:do-do-trung-thuc}. Ở phần mở đầu và phần công trình liên quan,
các tác giả dẫn một bộ đánh giá tên là M4 và ghi rằng bộ ấy cho thấy các chỉ
số trung thực thông dụng của attribution đặc trưng chỉ tương quan yếu với
nhau và có thể cho ra những bảng xếp hạng phương pháp mâu thuẫn
nhau~\autocite{p25userperception}. Sách không đọc bài M4 và không thêm khóa
trích dẫn cho nó, nên phát biểu này được thuật qua lời của bài báo hiện tại,
đúng theo cách chương~\ref{ch:chi-so-khong-do-do-trung-thuc} đã làm với các
công trình mà bài báo neo nhắc tới.

Mẩu bằng chứng ấy trả lời đúng một việc mà
mục~\ref{sec:ch09-doi-tuong-khac-nhau} đã giao. Trong ba thứ mà mục ấy nói là
đi qua được ranh giới, thứ ba là một phép kiểm mà phía attribution có thể chạy
ngay: đo mức đồng thuận giữa các chỉ số, việc không cần ground truth mà chỉ
cần chạy nhiều chỉ số trên cùng một tập. Phép kiểm ấy hóa ra đã được chạy, và
kết quả của nó là kết quả mà mục ấy dự đoán.

Nhưng phải đặt nó đúng chỗ, vì nó là thứ dễ bị đọc quá tay. Phép thử đồng
thuận giữa các chỉ số không phải phép đối chiếu với ground truth. Nhiều chỉ số
bất đồng với nhau thì không thể cùng đúng, nhưng nhiều chỉ số đồng thuận với
nhau cũng chưa chắc cái nào đúng, vì chúng có thể cùng sai theo một kiểu.
Chương~\ref{ch:chi-so-khong-do-do-trung-thuc} đã đọc đúng phép thử đồng thuận
ấy trên tám chỉ số của chuỗi suy luận và không rút từ nó một kết luận về độ
trung thực. Vậy điều đúng để nói là câu hỏi về mức đồng thuận nay đã được hỏi
cả với các chỉ số của attribution đặc trưng và câu trả lời cùng một hình dạng,
còn câu hỏi về đối chiếu với ground truth thì vẫn chưa được hỏi với chúng.
Phát biểu mà mục~\ref{sec:ch09-doi-tuong-khac-nhau} để ngỏ không đổi.

\index{đối chiếu chỉ số với người|)}%
